% This is text associated with the open textbook "Open Electromagnetics"
% See immediately below for author, date
% This work is licensed under the Creative Commons Attribution-ShareAlike 4.0 International License. (CC BY SA 4.0) 
% To view a copy of the license, visit http://creativecommons.org/licenses/by-sa/4.0/

%ORB TITLE Plane Waves in Lossy Regions
%ORB AUTHOR 0        # S.W. Ellingson, Virginia Tech (ellingson@vt.edu)
%ORB DATE 20191115   # yyyymmdd
%ORB ABSTRACT  

\label{m0130_Plane_Waves_Lossy_Regions} % <-- Must match name of this file and module directory name.  Don't mess with this!
\index{waves!plane|(}

The electromagnetic wave equations
\index{wave equation!electromagnetic}
for source-free regions consisting of possibly-lossy material are (see Section~\ref{m0128_Wave_Equations_Lossy_Regions}):
%
\begin{flalign}
\nabla^2\widetilde{\bf E} -\gamma^2 \widetilde{\bf E} &= 0 \label{m0130_eEg} \\
\nabla^2\widetilde{\bf H} -\gamma^2 \widetilde{\bf H} &= 0 \label{m0130_eHg}
\end{flalign} 
%
where
%
\begin{equation}
\gamma^2 \triangleq -\omega^2\mu\epsilon_c 
\label{m0130_egammadef}
\end{equation}
%
We now turn our attention to the question, what are the characteristics of waves that propagate in these conditions?
As in the lossless case, these equations permit waves having a variety of geometries including plane waves, cylindrical waves, and spherical waves.
In this section, we will consider the important special case of uniform plane waves.

To obtain the general expression for the uniform plane wave solution, we follow precisely the same procedure described in the section ``Uniform Plane Waves: Derivation.''\footnote{Depending on the version of this book, this section may appear in a different volume.}
Although that section presumed lossless media, the only difference in the present situation is that the real-valued constant $+\beta^2$ is replaced with the complex-valued constant $-\gamma^2$.  Thus, we obtain the desired solution through a simple modification of the solution for the lossless case.  For a wave exhibiting uniform magnitude and phase in planes of constant $z$, we find that the electric field is:
%
\begin{equation}
\widetilde{\bf E} = \hat{\bf x}\widetilde{E}_x + \hat{\bf y}\widetilde{E}_y
\end{equation}
%
where
%
\begin{flalign}
\widetilde{E}_x &= E_{x0}^+ e^{-\gamma z} + E_{x0}^- e^{+\gamma z} \label{m0130_eEx} \\
\widetilde{E}_y &= E_{y0}^+ e^{-\gamma z} + E_{y0}^- e^{+\gamma z} \label{m0130_eEy} 
\end{flalign}
%
where the complex-valued coefficients $E_{x0}^+$, $E_{x0}^-$, $E_{y0}^+$, and $E_{y0}^-$ are determined by boundary conditions (possibly including sources) outside the region of interest.  This result can be confirmed by verifying that Equations~\ref{m0130_eEx} and \ref{m0130_eEy} each satisfy Equation~\ref{m0130_eEg}.
Also, it may be helpful to note that these expressions are identical to those obtained for the voltage and current in lossy transmission lines, 
\index{transmission line!lossy}
\index{waves!transmission line analogy}
as described in the section ``Wave Equation for a TEM Transmission Line.''\footnote{Depending on the version of this book, this section may appear in a different volume.} 

Let's consider the special case of an $\hat{\bf x}$-polarized plane wave propagating in the $+\hat{\bf z}$ direction:
%
\begin{equation}
\widetilde{\bf E} = \hat{\bf x}E_{x0}^+ e^{-\gamma z}
\end{equation}
%
We established in Section~\ref{m0128_Wave_Equations_Lossy_Regions} that $\gamma$ may be written explicitly in terms of its real and imaginary components as follows:
%
\begin{equation}
\gamma = \alpha + j\beta
\end{equation}
%
where $\alpha$ and $\beta$ are positive real-valued constants depending on frequency ($\omega$) and constitutive properties of the medium; i.e., permittivity, permeability, and conductivity.  Thus:
%
\begin{flalign}
\widetilde{\bf E} &= \hat{\bf x}E_{x0}^+ e^{-\left(\alpha+j\beta\right) z} \nonumber \\
                  &= \hat{\bf x}E_{x0}^+ e^{-\alpha z} e^{-j\beta z}
\end{flalign}
%
Observe that the variation of phase with distance is determined by $\beta$ through the factor $e^{-j\beta z}$; thus, $\beta$ is the phase propagation constant 
\index{propagation constant!phase}
and plays precisely the same role as in the lossless case.
Observe also that the variation in magnitude is determined by $\alpha$ through the real-valued factor $e^{-\alpha z}$.  
Specifically, magnitude is reduced inverse-exponentially with increasing distance along the direction of propagation.  Thus, $\alpha$ is the attenuation constant
\index{propagation constant!attenuation}
and plays precisely the same role as the attenuation constant for a lossy transmission line.

%%%%%%%%%%%%%%%%%%%%%%%%%%%%%%%%%%%%%%%%%%%%%%%
%ORB HIGHLIGHT
\begin{mdframed}[backgroundcolor=yellow!20,nobreak=true] 
The presence of loss in material gives rise to a real-valued factor $e^{-\alpha z}$ which describes the attenuation of the wave with propagation (in this case, along $z$) in the material.
\end{mdframed}
%%%%%%%%%%%%%%%%%%%%%%%%%%%%%%%%%%%%%%%%%%%%%%%

We may continue to exploit the similarity of the potentially-lossy and lossless plane wave results to quickly ascertain the characteristics of the magnetic field.  In particular, the plane wave relationships 
\index{plane wave relationships}
apply exactly as they do in the lossless case.  These relationships are:
%
%
\begin{equation}
\widetilde{\bf H} = \frac{1}{\eta} \hat{\bf k} \times \widetilde{\bf E} 
\label{m0130_ePWRH}
\end{equation}
%
\begin{equation}
\widetilde{\bf E} = -\eta \hat{\bf k} \times \widetilde{\bf H} 
\label{m0130_ePWRE}
\end{equation}
%
where $\hat{\bf k}$ is the direction of propagation and $\eta$ is the wave impedance.
\index{wave impedance}
In the lossless case, $\eta=\sqrt{\mu/\epsilon}$; however, in the possibly-lossy case we must replace $\epsilon=\epsilon'$ with $\epsilon_c=\epsilon'-j\epsilon''$.  Thus:
%
\begin{flalign}
\eta \rightarrow \eta_c &= \sqrt{\frac{\mu}{\epsilon_c}} = \sqrt{\frac{\mu}{\epsilon'-j\epsilon''}} \nonumber \\
                        &= \sqrt{\frac{\mu}{\epsilon'}}\sqrt{\frac{1}{1-j\left(\epsilon''/\epsilon'\right)}} 
\end{flalign}
%
Thus:
%
\begin{equation}
\boxed{
\eta_c = \sqrt{\frac{\mu}{\epsilon'}} \cdot \left[ 1-j\frac{\epsilon''}{\epsilon'} \right]^{-1/2}
}
\label{m0130_eetac}
\end{equation}
%
Remarkably, we find that the wave impedance for a lossy material is equal to $\sqrt{\mu/\epsilon}$ -- the wave impedance we would calculate if we neglected loss (i.e., assumed $\sigma=0$) -- times a correction factor that accounts for the loss.  This correction factor is complex-valued; therefore, ${\bf E}$ and ${\bf H}$ are not in phase when propagating through lossy material.  We now see that in the phasor domain:
%
\begin{equation}
\widetilde{\bf H} = \frac{1}{\eta_c} \hat{\bf k} \times \widetilde{\bf E} \label{m0130_eEpw}
\end{equation}
%
\begin{equation}
\widetilde{\bf E} = -\eta_c \hat{\bf k} \times \widetilde{\bf H} \label{m0130_eHpw}
\end{equation}   

%%%%%%%%%%%%%%%%%%%%%%%%%%%%%%%%%%%%%%%%%%%%%%%
%ORB HIGHLIGHT
\begin{mdframed}[backgroundcolor=yellow!20,nobreak=true] 
The plane wave relationships in media which are possibly lossy are given by Equations~\ref{m0130_eEpw} and \ref{m0130_eHpw}, with the complex-valued wave impedance given by Equation~\ref{m0130_eetac}.
\end{mdframed}
%%%%%%%%%%%%%%%%%%%%%%%%%%%%%%%%%%%%%%%%%%%%%%%

\index{waves!plane|)}

